\begin{table}[htbp]
\centering
\small
\caption{Bloco 2 (google/gemini-2.5-flash-lite): viés de ordem dos exemplos.}
\label{tab:bloco2_E1_ordem_exemplos_flashlite}
\begin{tabular}{lrrrrrr}
\toprule
 & \multicolumn{3}{c}{Acur.\ vs espec.\ (\%)} & \multicolumn{3}{c}{$\kappa$} \\ \cmidrule(lr){2-4} \cmidrule(lr){5-7}
Ordenação & n = 4 & n = 10 & n = 20 & n = 4 & n = 10 & n = 20 \\
\midrule
classe 0 primeiro & 85,4 & 89,3 & 91,8 & 0,707 & 0,785 & 0,835 \\
classe 1 primeiro & 79,5 & 84,5 & 82,3 & 0,592 & 0,692 & 0,649 \\
embaralhada & 89,7 & 89,0 & 90,8 & 0,795 & 0,781 & 0,815 \\
alternada & 84,5 & 90,7 & 92,6 & 0,690 & 0,814 & 0,851 \\
\midrule
\textbf{amplitude (máx $-$ mín)} & \textbf{10,3} & \textbf{6,2} & \textbf{10,3} &  &  &  \\
\bottomrule
\multicolumn{7}{p{0.95\linewidth}}{\footnotesize Mesmos exemplos (estratégia mista, especialista E1) em quatro ordenações. Acur.\ vs espec.\ = acurácia (LLM vs especialista): fração dos pontos de teste em que o LLM reproduz o rótulo do especialista. $\kappa$ = Cohen's $\kappa$ (faixas de Landis e Koch, 1977). Médias das sementes.} \\
\end{tabular}
\end{table}
